% =====================================================================
%  Reunión de tesis con Prof. Enrique Espina — lunes 23-11-2026, 17:00
%  Felipe Gajardo Ávalos · Magíster en Ciencias de la Ingeniería, USACH
%  Tema: margen de error del terciario y lazo rápido (acuerdo 3, acta 16-11)
%  v2 (19-11): rehecha tras revisar el PDF compilado (sumatorias grandes,
%  recuadros desbordados, diagrama del lazo confuso, exceso de texto).
%   L1  Niveles de control (actualizado)
%   L2  Dos tipos de desbalance: energía y margen
%   L3  El lazo rápido: plan + reacción proporcional (política afín, alpha_k)
%   L4  Cuánto margen guardar: el modelo del error
%   L5  El margen de energía crece con el horizonte
%   L6  Terciario: el LP completo
%   L7  Lazo rápido coherente con el terciario
%   L8  Decisiones que necesito validar
%  Compilar con XeLaTeX (recomendado) o pdfLaTeX, dos pasadas.
%  Requiere fig/fondo_usach_contenido.png
% =====================================================================
\PassOptionsToPackage{table}{xcolor}
\documentclass[aspectratio=169,10pt,t]{beamer}

\usefonttheme{professionalfonts}
% XeLaTeX/LuaLaTeX: fontspec + polyglossia. pdfLaTeX: inputenc + babel.
\usepackage{iftex}
\ifPDFTeX
  \usepackage[utf8]{inputenc}
  \usepackage[T1]{fontenc}
  \usepackage{lmodern}
  \usepackage[spanish,es-noshorthands,es-nodecimaldot]{babel}
\else
  \usepackage{fontspec}
  \usepackage{polyglossia}
  \setmainlanguage{spanish}
\fi
\usepackage{amsmath,amssymb,bm}
\usepackage{exscale}   % sumatorias del tamaño de la letra (sin esto salen gigantes)
\usepackage{graphicx}
\usepackage{array}
\usepackage{tikz}
\usetikzlibrary{arrows.meta,positioning,calc,fit,backgrounds,shapes.misc,decorations.pathreplacing}
\graphicspath{{fig/}}
\hyphenpenalty=10000\exhyphenpenalty=10000   % sin guiones de corte en los recuadros

% ------------------------------- Colores USACH -----------------------
\definecolor{usachT}{HTML}{009688}
\definecolor{usachTd}{HTML}{00695C}
\definecolor{usachO}{HTML}{E77500}
\definecolor{usachOl}{HTML}{FDF0E3}
\definecolor{usachR}{HTML}{C62828}
\definecolor{verdeL}{HTML}{2E7D32}
\definecolor{azulP}{HTML}{1F5FAF}
\definecolor{grisM}{HTML}{55606E}
\colorlet{cT}{usachO}\colorlet{cL}{verdeL}\colorlet{cS}{usachT}\colorlet{cP}{azulP}

\setbeamercolor{normal text}{fg=black}
\setbeamercolor{structure}{fg=usachT}
\setbeamercolor{frametitle}{fg=usachT}
\setbeamertemplate{navigation symbols}{}
\setbeamerfont{frametitle}{size=\Large}
\setbeamersize{text margin left=2.75cm,text margin right=0.55cm}
\usebackgroundtemplate{\includegraphics[width=\paperwidth,height=\paperheight]{fondo_usach_contenido.png}}
\setbeamertemplate{headline}{%
  \begin{tikzpicture}[remember picture,overlay]
    \node[fill=usachT!55,text=white,rounded corners=1.5pt,font=\tiny,inner sep=2pt,anchor=north east]
      at ([xshift=-0.22cm,yshift=-0.16cm]current page.north east) {\insertframenumber/\inserttotalframenumber};
  \end{tikzpicture}}
\setbeamertemplate{frametitle}{\vskip0.36cm\usebeamerfont{frametitle}\usebeamercolor[fg]{frametitle}\insertframetitle\par\vskip0.04cm}
\setbeamertemplate{footline}{\vbox to 0.62cm{}}

\tikzset{
  tag/.style={fill=#1,text=white,rounded corners=3pt,align=center,text width=2.3cm,minimum height=1.95cm,inner sep=3pt,font=\footnotesize},
  ent/.style={draw=gray!55,fill=white,rounded corners=2pt,align=left,text width=4.55cm,minimum height=1.95cm,inner sep=4pt,font=\footnotesize,line width=.6pt},
  dec/.style={draw=#1,line width=1.3pt,fill=#1!10,rounded corners=3pt,align=center,text width=2.95cm,minimum height=1.95cm,inner sep=3pt,font=\footnotesize},
  fin/.style={-{Latex[length=2.4mm]},line width=1.1pt,gray!60!black},
  hdr/.style={font=\footnotesize\bfseries,text=grisM},
  out/.style={draw=#1,line width=.9pt,fill=white,rounded corners=2pt,align=left,text width=2.95cm,minimum height=0.8cm,inner sep=3pt,font=\footnotesize},
  dest/.style={fill=#1,text=white,rounded corners=5pt,align=center,text width=1.7cm,minimum height=0.42cm,inner sep=2pt,font=\footnotesize\bfseries},
  fo/.style={-{Latex[length=2.2mm]},line width=1.2pt,#1},
  jd/.style={circle,fill=#1,inner sep=0pt,minimum size=3.2pt},
  kc/.style={draw=#1,line width=.8pt,rounded corners=2pt,fill=white,align=left,inner sep=4pt,font=\fontsize{7}{8.6}\selectfont},
  blq/.style={draw=#1,line width=1.1pt,fill=white,rounded corners=2pt,align=center,font=\small,inner sep=3pt,minimum height=0.95cm},
  sm/.style={circle,draw=#1,line width=1.1pt,fill=white,inner sep=0pt,minimum size=0.55cm,font=\small\bfseries},
  ar/.style={-{Latex[length=2.3mm]},line width=1.2pt,#1},
}
\newcommand{\kt}[2]{{\fontsize{7.6}{9}\selectfont\bfseries\color{#1}#2}}

\begin{document}

% =====================================================================
%  L1 — NIVELES
% =====================================================================
\begin{frame}{Niveles de control: qué mide, decide e informa}
  \vspace{-0.28cm}
  \centering
  \resizebox{!}{6.1cm}{%
  \begin{tikzpicture}[x=1cm,y=1cm]
    \node[hdr] at (1.3,9.55) {NIVEL};
    \node[hdr] at (5.2,9.55) {1.\; MIDE O RECIBE};
    \node[hdr] at (9.55,9.55) {2.\; DECIDE};
    \node[hdr] at (13.25,9.55) {3.\; INFORMA};
    \node[hdr] at (16.35,9.55) {4.\; A QUIÉN};
    \foreach \y in {7.2,5.0,2.8}{\draw[gray!30,line width=.5pt] (0,\y) -- (17.3,\y);}

    % TERCIARIO
    \node[tag=cT] (gT) at (1.3,8.2) {\textbf{TERCIARIO}\\[1pt] uno, central\\ cada 15 min\\[3pt] \colorbox{white}{\scriptsize\bfseries\textcolor{cT}{PC host}}};
    \node[ent] (eT) at (5.2,8.2) {%
      $\mathcal E_k,\ E_k$: envolvente y energía de cada MR\\
      $\hat D_\phi$: desbalance previsto por fase\\
      $\sigma_\phi,\ a_\phi$: estadística de su error\\
      $c(t)$: tarifa\\[2pt]
      \textcolor{cS}{de: secundarios (vía su IC) e historia}};
    \node[dec=cT] (T) at (9.55,8.2) {\textbf{Solo energía}\\[1pt]
      mín.\ costo de compra\\ LP con $H=12$ h\\ aparta energía\\ y margen};
    \node[out=cT] (oT1) at (13.25,8.67) {$P^*_k,\ \alpha_k$\\ {\scriptsize plan y participación}};
    \node[out=cT] (oT2) at (13.25,7.73) {$b_\phi,\ \rho_\phi,\ R^{\varepsilon}_\phi$\\ {\scriptsize banda de la fase}};
    \node[dest=cP] (dT1) at (16.35,8.67) {IC$_k$};
    \node[dest=cL] (dT3) at (16.35,7.73) {Lazo rápido};
    \node[font=\scriptsize,text=grisM] at (16.35,7.3) {(uno por fase)};

    % LAZO RÁPIDO
    \node[tag=cL] (gL) at (1.3,6.1) {\textbf{LAZO RÁPIDO}\\[1pt] uno por fase\\ cada $\sim$1 s\\[3pt] \colorbox{white}{\scriptsize\bfseries\textcolor{cL}{controlador}}};
    \node[ent] (eL) at (5.2,6.1) {%
      $P_{h,a},\ P_{h,b},\ P_{h,c}$: potencia por fase\\
      $b_\phi,\ \rho_\phi,\ R^{\varepsilon}_\phi$: banda de la fase\\[2pt]
      \textcolor{cT}{de: cabecera y terciario}};
    \node[dec=cL] (L) at (9.55,6.1) {\textbf{Solo balance}\\[1pt]
      iguala las 3 fases\\ $e_\phi=P_{h,\phi}-\bar P_h$\\ PI en $b_\phi\pm R^{\varepsilon}_\phi$};
    \node[out=cL] (oL1) at (13.25,6.35) {$\Delta P_\phi$\\ {\scriptsize corrección de la fase}};
    \node[dest=cP] (dL1) at (16.35,6.35) {IC de la fase $\phi$};
    \node[font=\scriptsize,text=usachR!85!black,align=center] at (14.8,5.5) {no habla con secundarios ni DER};

    % SECUNDARIO
    \node[tag=cS] (gS) at (1.3,3.9) {\textbf{SECUNDARIO}\\[1pt] uno por MR\\ segundos\\[3pt] \colorbox{white}{\scriptsize\bfseries\textcolor{cS}{controlador}}};
    \node[ent] (eS) at (5.2,3.9) {%
      $V_{dc,k}$: tensión de su barra DC\\
      SoC: de su BESS y su EV\\
      $P_{\mathrm{PV},k}$: FV disponible\\[2pt]
      \textcolor{cS}{de: medidas locales de la MR}};
    \node[dec=cS] (S) at (9.55,3.9) {Restaura $V_{dc,k}$\\ reparte entre sus DER\\ cuida límites BESS/EV\\ calcula su envolvente\\
      \textcolor{usachR}{no ve la cabecera}};
    \node[out=cS] (oS1) at (13.25,4.37) {$V^*_{dc,d}$\\ {\scriptsize referencia de cada DER $d$}};
    \node[out=cS] (oS2) at (13.25,3.43) {$\mathcal E_k,\;E_k$\\ {\scriptsize envolvente ya ajustada}};
    \node[dest=cP] (dS1) at (16.35,4.37) {Primario DER};
    \node[dest=cT] (dS2) at (16.35,3.43) {Terciario};
    \node[font=\scriptsize,text=grisM] at (16.35,3.0) {(vía su IC)};

    % PRIMARIO
    \node[tag=cP] (gP) at (1.3,1.7) {\textbf{PRIMARIO}\\[1pt] uno por\\ convertidor\\ milisegundos\\[3pt] \colorbox{white}{\scriptsize\bfseries\textcolor{cP}{controlador}}};
    \node[ent] (eP) at (5.2,1.7) {%
      DER: $V_{dc,k}$ y su referencia\\
      IC$_k$: $P_{\mathrm{ref},k}=P^*_k+\alpha_k\,\Delta P_\phi$\\[2pt]
      \textcolor{cP}{de: secundario, terciario y lazo}};
    \node[dec=cP] (P) at (9.55,1.7) {\textbf{DER}: droop $V_{dc}$--$P$\\[3pt]
      \textbf{IC$_k$}: grid-following\\ sigue $P_{\mathrm{ref},k}$\\ ($Q_{\mathrm{ref},k}=0$)};
    \node[out=cP] (oP1) at (13.25,2.17) {DER: $i_d$\\ {\scriptsize corriente a la barra DC}};
    \node[out=cP] (oP2) at (13.25,1.23) {IC$_k$: $p_k$\\ {\scriptsize potencia inyectada}};
    \node[dest=grisM] (dP1) at (16.35,2.17) {Barra DC};
    \node[dest=grisM] (dP2) at (16.35,1.23) {Fase $\phi$};
    \node[font=\scriptsize,text=grisM] at (16.35,0.7) {(la mide la cabecera)};

    \foreach \n in {T,L,S,P}{\draw[fin] (e\n.east) -- (\n.west);}
    \foreach \n/\c/\yc in {T/cT/8.2,S/cS/3.9,P/cP/1.7}{
      \draw[line width=1.2pt,\c] (\n.east) -- (11.4,\yc);
      \draw[fo=\c] (11.4,\yc) |- (o\n1.west);
      \draw[fo=\c] (11.4,\yc) |- (o\n2.west);
      \node[jd=\c] at (11.4,\yc) {};
    }
    \draw[fo=cL] (L.east) -- (11.4,6.1) |- (oL1.west);
    \draw[fo=cT] (oT1.east) -- (dT1.west);
    \draw[fo=cT] (oT2.east) -- (dT3.west);
    \draw[fo=cL] (oL1.east) -- (dL1.west);
    \draw[fo=cS] (oS1.east) -- (dS1.west);
    \draw[fo=cS] (oS2.east) -- (dS2.west);
    \draw[fo=cP] (oP1.east) -- (dP1.west);
    \draw[fo=cP] (oP2.east) -- (dP2.west);
  \end{tikzpicture}}
\end{frame}

% =====================================================================
%  L2 — DOS TIPOS DE DESBALANCE
% =====================================================================
\begin{frame}{Dos tipos de desbalance: energía y margen}
  \vspace{-0.30cm}
  \begin{columns}[T,onlytextwidth]
    \column{.52\textwidth}
    \centering
    \begin{tikzpicture}[x=0.86cm,y=0.82cm,font=\fontsize{6.6}{7.8}\selectfont]
      \node[font=\fontsize{7.4}{9}\selectfont\bfseries,anchor=west] at (-0.3,4.15) {Fase $a$: cuánto se aparta del promedio [kW]};
      \fill[cP!10] (1,1.5) rectangle (5,2.87);
      \fill[cT!20] (1,0) rectangle (5,1.5);
      \draw[gray!60,-{Latex[length=1.5mm]}] (0,0) -- (6.4,0);
      \draw[gray!60,-{Latex[length=1.5mm]}] (0,0) -- (0,3.6);
      \foreach \x/\l in {0/17,1/18,2/19,3/20,4/21,5/22,6/23}{\draw[gray!60] (\x,0) -- (\x,-0.08) node[below]{\l};}
      \node at (3,-0.62) {hora};
      \foreach \v in {1,2,3}{\draw[gray!60] (0,\v) -- (-0.08,\v) node[left]{\v};}
      \draw[cT,line width=1.4pt] (0,0.02) -- (1,0.02) -- (1,1.5) -- (5,1.5) -- (5,0.02) -- (6,0.02);
      \draw[cP,line width=.8pt,densely dashed] (1,2.87) -- (5,2.87);
      \node[cT!80!black,align=center,font=\fontsize{6.6}{7.8}\selectfont\bfseries] at (3,0.72) {previsible\\ $1{,}5$ kW $\times$ 4 h $=6$ kWh};
      \node[cP,anchor=south,font=\fontsize{6.6}{7.8}\selectfont\bfseries] at (3,2.9) {margen para el error: $+1{,}37$ kW};
      \draw[usachR,line width=1pt] plot[smooth] coordinates {(0,0.18) (0.5,0.08) (0.92,0.25) (1.15,1.45) (1.5,1.85) (2,2.05) (2.5,1.85) (3,2.12) (3.5,1.92) (4,2.15) (4.5,1.82) (4.9,1.92) (5.12,0.45) (5.5,0.18) (6,0.28)};
      \node[usachR,anchor=west,align=left,font=\fontsize{6.6}{7.8}\selectfont\bfseries] at (5.15,1.25) {lo que\\ ocurre};
    \end{tikzpicture}\\[4pt]
    \raggedright
    \begin{tikzpicture}\node[kc=usachR,text width=6.0cm]{%
      \kt{usachR}{Formulación anterior}\\[2pt]
      Todo el desbalance se cubría con margen de potencia de 15 min. El lazo rápido gastaba los 6 kWh previsibles sin que nadie los hubiera apartado: la BESS se agotaba a las 20:40.};\end{tikzpicture}

    \column{.46\textwidth}
    \begin{tikzpicture}\node[kc=cT,text width=5.55cm]{%
      \kt{cT!85!black}{Lo previsible: $\hat e_\phi$}\\[2pt]
      Clientes monofásicos de siempre, todas las tardes. Dura horas.\\[2pt]
      $\Rightarrow$ se aparta como \textbf{energía comprometida} en la trayectoria de SoC de cada MR.};\end{tikzpicture}\\[5pt]
    \begin{tikzpicture}\node[kc=cP,text width=5.55cm]{%
      \kt{cP}{El error del pronóstico: $\xi_\phi$}\\[2pt]
      Lo que no se sabe de antemano. Media cero, pero suele persistir varias horas.\\[2pt]
      $\Rightarrow$ se aparta como \textbf{margen}: de potencia ($r_k$) y de energía ($Z_\phi$).};\end{tikzpicture}\\[5pt]
    \begin{tikzpicture}\node[kc=cL,text width=5.55cm]{%
      \kt{verdeL}{Los roles no cambian}\\[2pt]
      El lazo rápido es el único que balancea. El terciario solo aparta lo que el lazo va a gastar; su objetivo sigue siendo solo energía.};\end{tikzpicture}
  \end{columns}
\end{frame}

% =====================================================================
%  L3 — POLÍTICA AFÍN Y FACTORES DE PARTICIPACIÓN
% =====================================================================
\begin{frame}{El lazo rápido: plan más reacción proporcional}
  \vspace{-0.30cm}
  \begin{columns}[T,onlytextwidth]
    \column{.47\textwidth}
    \begin{tikzpicture}\node[kc=cS,text width=5.65cm]{%
      \kt{usachTd}{Lo que hace cada IC de la fase $\phi$}\\[4pt]
      \hspace*{4pt}$P_{\mathrm{IC},k}(t)=\underbrace{P^*_k(t)}_{\text{plan}}+\underbrace{\alpha_k}_{\text{su parte}}\,\underbrace{\Delta P_\phi(t)}_{\text{corrección de la fase}}$\\[4pt]
      $\alpha_k\ge0$, $\ \sum_{k\in\mathcal N_\phi}\alpha_k=1$: \textbf{factores de participación}.\\[2pt]
      Plan fijo más reacción lineal al error medido: es una \textbf{política afín de reserva} con respuesta diagonal, la estructura de la reserva secundaria clásica [Warrington et al., 2013].};\end{tikzpicture}\\[5pt]
    \begin{tikzpicture}\node[kc=cT,text width=5.65cm]{%
      \kt{cT!85!black}{Consecuencia: el terciario sabe qué va a gastar cada MR}\\[3pt]
      \hspace*{4pt}$\hat y_k=\alpha_k\,b_\phi$\quad energía comprometida [kW]\\[2pt]
      \hspace*{4pt}$r_k=\alpha_k\,R^{\varepsilon}_\phi$\quad margen para el error [kW]\\[3pt]
      Criterio propuesto: $\alpha_k=E_k/\sum_j E_j$ $\Rightarrow$ todas las MR de la fase bajan \textbf{el mismo $\Delta$SoC}.};\end{tikzpicture}

    \column{.51\textwidth}
    \begin{tikzpicture}\node[kc=grisM,text width=6.15cm]{%
      \kt{grisM}{Por qué el reparto planificado debe ser el del lazo}\\[2pt]
      Fase $a$, evento previsto de $1{,}5$ kW durante 4 h.\\[3pt]
      \begin{tabular}{@{}l@{\hspace{6pt}}c@{\hspace{6pt}}c@{\hspace{6pt}}c@{}}
        & MR$_1$ & MR$_2$ & MR$_3$\\\hline
        $E_k$ [kWh] & 10 & 10 & 5\\[2pt]
        \multicolumn{4}{@{}l}{\textcolor{usachR}{\textbf{Reparto libre en el terciario}}}\\
        plan [kW] & 1,5 & 0 & 0\\
        lo que hace el lazo [kW] & 0,6 & 0,6 & 0,3\\[2pt]
        \multicolumn{4}{@{}l}{\textcolor{verdeL}{\textbf{Con $\alpha_k$ fijo $=0{,}4/0{,}4/0{,}2$}}}\\
        plan $=$ lo que hace el lazo [kW] & 0,6 & 0,6 & 0,3\\
        energía en 4 h [kWh] & 2,4 & 2,4 & 1,2\\
        $\Delta$SoC & 24\,\% & 24\,\% & 24\,\%\\
      \end{tabular}\\[4pt]
      Con reparto libre, MR$_2$ y MR$_3$ gastan energía que nadie les apartó. Con $\alpha_k$ fijo, lo planificado es lo que ocurre y todo sigue siendo un \textbf{LP}. Optimizar $\alpha_k$ lo vuelve bilineal: queda como extensión.};\end{tikzpicture}
  \end{columns}
\end{frame}

% =====================================================================
%  L4 — MODELO DEL ERROR
% =====================================================================
\begin{frame}{Cuánto margen guardar: el modelo del error}
  \vspace{-0.30cm}
  \centering
  \begin{tikzpicture}[x=1cm,y=1cm]
    \node[kc=cS,text width=3.6cm,minimum height=1.75cm,anchor=north west] (A) at (0,0) {%
      \kt{usachTd}{1. Pronóstico simple}\\[2pt]
      mismo tramo de 15 min de los últimos $n$ días:\\[2pt]
      $\hat D_\phi(t)=\frac1n\sum_{d=1}^{n}D_\phi(t-96d)$};
    \node[kc=cS,text width=3.6cm,minimum height=1.75cm,anchor=north west] (B) at (4.25,0) {%
      \kt{usachTd}{2. Su error, en la historia}\\[2pt]
      $\xi_\phi=D_\phi-\hat D_\phi$ [kW]\\[1pt]
      $\sigma_\phi(h)$: tamaño típico, por hora\\
      $a_\phi$: cuánto persiste de un tramo al siguiente};
    \node[kc=cT,text width=3.6cm,minimum height=1.75cm,anchor=north west] (C) at (8.5,0) {%
      \kt{cT!85!black}{3. Margen al riesgo $\varepsilon$}\\[2pt]
      $R^{\varepsilon}_\phi=q_{1-\varepsilon}(\xi_\phi)+m$\\[2pt]
      alcanza en una fracción $1-\varepsilon$ de los casos};
    \draw[ar=grisM] (A.east) -- (B.west);
    \draw[ar=grisM] (B.east) -- (C.west);
  \end{tikzpicture}\\[6pt]
  \raggedright
  \begin{columns}[T,onlytextwidth]
    \column{.40\textwidth}
    \centering
    \begin{tikzpicture}[x=0.62cm,y=0.85cm,font=\fontsize{6.6}{7.8}\selectfont]
      \fill[usachR!25] (2.33,0) -- plot[domain=2.33:3.4,samples=20] (\x,{1.5*exp(-((\x)^2)/2)}) -- (3.4,0) -- cycle;
      \draw[gray!60,-{Latex[length=1.4mm]}] (-3.6,0) -- (3.9,0) node[below left]{$\xi_\phi$};
      \draw[cT,line width=1pt] plot[domain=-3.4:3.4,samples=60,smooth] (\x,{1.5*exp(-((\x)^2)/2)});
      \draw[usachR,line width=.8pt] (2.33,0) -- (2.33,1.0);
      \node[usachR,anchor=south] at (2.33,1.0) {$q_{1-\varepsilon}$};
      \node[usachR,anchor=west] at (2.75,0.32) {área $\varepsilon$};
      \node at (0,-0.25) {0};
    \end{tikzpicture}\\[3pt]
    \raggedright
    \begin{tikzpicture}\node[kc=grisM,text width=4.65cm]{%
      \kt{grisM}{Qué cubre $m$}\\[2pt]
      Lo que el promedio de 15 min no ve: variación dentro del tramo, error de medida y de seguimiento del PI.};\end{tikzpicture}

    \column{.58\textwidth}
    \begin{tikzpicture}\node[kc=cT,text width=7.0cm]{%
      \kt{cT!85!black}{Cómo se obtiene $q_{1-\varepsilon}$ (con $\varepsilon=1\,\%$)}\\[3pt]
      \begin{tabular}{@{}p{4.5cm}@{\hspace{6pt}}l@{}}
        Gaussiana: $z\,\sigma_\phi$, $z=\Phi^{-1}(1-\varepsilon)$ & $2{,}33\,\sigma_\phi$\\
        Sin supuesto (Cantelli) & $9{,}95\,\sigma_\phi$\\
        \textbf{Cuantil empírico de la historia} & \textbf{propuesta}\\
      \end{tabular}\\[4pt]
      Precedente: la reserva FRR europea cubre el \textbf{99\,\% de los desbalances históricos}, por separado hacia arriba y hacia abajo [SOGL, art.~157].\\[4pt]
      \kt{grisM}{Ejemplo (supuestos)}\\[1pt]
      $\sigma_a=0{,}5$ kW, $z=2{,}33$, $m=0{,}2$ kW $\ \Rightarrow\ R^{\varepsilon}_a=1{,}37$ kW\\
      Repartido con $\alpha=0{,}4/0{,}4/0{,}2$: $\ r_k=0{,}55/0{,}55/0{,}27$ kW};\end{tikzpicture}
  \end{columns}
\end{frame}

% =====================================================================
%  L5 — ENERGÍA EN LA TRAYECTORIA
% =====================================================================
\begin{frame}{El margen de energía crece con el horizonte}
  \vspace{-0.30cm}
  \begin{columns}[T,onlytextwidth]
    \column{.49\textwidth}
    \begin{tikzpicture}\node[kc=cS,text width=5.9cm]{%
      \kt{usachTd}{Energía que entrega el IC$_k$ hasta el tramo $t$}\\[3pt]
      \hspace*{4pt}$W_k(t)=\sum_{s\le t}\tilde P_k(s)\,\Delta t\ +\ \alpha_k\sum_{s\le t}\xi_\phi(s)\,\Delta t$\\[1pt]
      \hspace*{4pt}{\color{grisM}\scriptsize plan $+$ previsto \hspace{2.1cm} error acumulado}\\[4pt]
      \kt{usachTd}{El error acumulado crece (modelo AR(1), persistencia $a$)}\\[3pt]
      \hspace*{4pt}$\sigma_{E,\phi}(n)=\sigma_\phi\,\Delta t\,\sqrt{f(n)}$\quad tras $n$ tramos\\[3pt]
      \hspace*{4pt}$f(n)=\dfrac{n(1+a)}{1-a}-\dfrac{2a(1-a^{n})}{(1-a)^2}$\\[3pt]
      $a=0$: $f=n$ (independientes)\quad $a\to1$: $f\to n^2$ (persistente)};\end{tikzpicture}\\[5pt]
    \begin{tikzpicture}\node[kc=cT,text width=5.9cm]{%
      \kt{cT!85!black}{Restricción (c): toda la trayectoria}\\[3pt]
      \hspace*{4pt}$W^{\min}_k(t)+\alpha_kZ^{\downarrow}_\phi(t)\ \le\ \sum_{s\le t}\tilde P_k(s)\Delta t$\\[2pt]
      \hspace*{4pt}$\sum_{s\le t}\tilde P_k(s)\Delta t\ \le\ W^{\max}_k(t)-\alpha_kZ^{\uparrow}_\phi(t)$\\[3pt]
      \hspace*{4pt}$Z_\phi(t)=\big(z\,\sigma_\phi\sqrt{f(t+1)}+m\big)\,\Delta t$\\[3pt]
      Primer tramo: $Z_\phi=R^{\varepsilon}_\phi\,\Delta t$. El antiguo $\tau$ \textbf{es} $\Delta t=15$ min, porque cada 15 min se re-planifica con el SoC medido.};\end{tikzpicture}

    \column{.49\textwidth}
    \centering
    \begin{tikzpicture}[x=1.3cm,y=0.92cm,font=\fontsize{6.6}{7.8}\selectfont]
      \node[font=\fontsize{7.4}{9}\selectfont\bfseries,anchor=west] at (-0.4,3.75) {Margen de energía de la fase, $Z_\phi$ [kWh]};
      \draw[gray!60,-{Latex[length=1.5mm]}] (0,0) -- (4.3,0);
      \draw[gray!60,-{Latex[length=1.5mm]}] (0,0) -- (0,3.45);
      \foreach \x in {0,1,2,3,4}{\draw[gray!60] (\x,0) -- (\x,-0.07) node[below]{\x};}
      \node at (2,-0.62) {horizonte [h]};
      \foreach \v in {1,2,3}{\draw[gray!60] (0,\v) -- (-0.07,\v) node[left]{\v}; \draw[gray!15] (0,\v) -- (4.1,\v);}
      \draw[usachR,line width=1.3pt,densely dashed] (0.25,0.3425) -- (4,0.3425);
      \draw[cP,line width=1.3pt] plot[smooth] coordinates {(0.25,0.341) (0.5,0.462) (0.75,0.554) (1,0.633) (1.5,0.763) (2,0.874) (3,1.059) (4,1.215)};
      \draw[cT,line width=1.3pt] plot[smooth] coordinates {(0.25,0.341) (0.5,0.603) (0.75,0.847) (1,1.075) (1.5,1.491) (2,1.862) (3,2.500) (4,3.036)};
      \fill[cT] (4,3.036) circle (1.6pt);
      \node[cT!85!black,anchor=south east] at (3.95,3.06) {3,04};
      \node[cP,anchor=south east] at (3.95,1.24) {1,22};
    \end{tikzpicture}\\[2pt]
    \raggedright
    {\fontsize{6.8}{8.2}\selectfont
    \textcolor{cT}{\rule[2pt]{0.5cm}{1.3pt}} error persistente ($a=0{,}8$)\quad
    \textcolor{cP}{\rule[2pt]{0.5cm}{1.3pt}} independiente ($a=0$)\\
    \textcolor{usachR}{\rule[2pt]{0.5cm}{1.3pt}} $R^{\varepsilon}\tau$ constante (formulación anterior)\\[1pt]
    \textcolor{grisM}{$\sigma=0{,}5$ kW, $z=2{,}33$, $m=0{,}2$ kW (supuestos)}\par}
    \vspace{4pt}
    \begin{tikzpicture}\node[kc=grisM,text width=5.9cm]{%
      MR$_1$ ($\alpha_1=0{,}4$), 4 h, $a=0{,}8$: comprometida $2{,}4$ kWh $+$ margen $0{,}4\times3{,}04=1{,}21$ kWh $=3{,}6$ kWh de 10. Con $R^{\varepsilon}\tau$: solo $0{,}14$ kWh.};\end{tikzpicture}
  \end{columns}
\end{frame}

% =====================================================================
%  L6 — LP COMPLETO
% =====================================================================
\begin{frame}{Terciario: el LP completo}
  \vspace{-0.30cm}
  \begin{columns}[T,onlytextwidth]
    \column{.38\textwidth}
    \begin{tikzpicture}\node[kc=cS,text width=4.5cm]{%
      \kt{usachTd}{Variables de decisión}\\[2pt]
      $P^*_k(t)$ [kW]: plan de cada IC\\
      $\rho_\phi(t)$ [kW]: residual aceptado\\
      $g(t)$ [kW]: compra a la S/E};\end{tikzpicture}\\[5pt]
    \begin{tikzpicture}\node[kc=grisM,text width=4.5cm]{%
      \kt{grisM}{Parámetros que entran}\\[2pt]
      $c(t)$: tarifa\\
      $\hat P_{L,\mathrm{tot}}(t)$, $\hat D_\phi(t)$: pronósticos\\
      $P^{\min/\max}_k$, $W^{\min/\max}_k$: envolvente\\
      $\alpha_k$: participación\\
      $R^{\varepsilon}_\phi(t)$, $Z_\phi(t)$: margen\\
      $\bar\rho$: residual admisible\\
      $W^{\mathrm{fin}}_k$: energía al final del horizonte};\end{tikzpicture}\\[5pt]
    \begin{tikzpicture}\node[kc=cP,text width=4.5cm]{%
      \kt{cP}{Cada 15 min (horizonte rodante)}\\[2pt]
      medir SoC $\to$ resolver con $T=48$ tramos de $0{,}25$ h $\to$ aplicar solo el primero};\end{tikzpicture}

    \column{.60\textwidth}
    \begin{tikzpicture}\node[kc=cT,text width=7.2cm]{%
      \kt{cT!85!black}{Objetivo: solo energía}\\[2pt]
      \hspace*{4pt}$\min\ \sum_{t}c(t)\,g(t)\,\Delta t$\\[4pt]
      \kt{cT!85!black}{Definiciones}\\[2pt]
      \hspace*{4pt}$S_\phi=\sum_{k\in\mathcal N_\phi}P^*_k,\quad \bar S=\tfrac13\sum_\psi S_\psi,\quad \hat e_\phi=\hat D_\phi-(S_\phi-\bar S)$\\[2pt]
      \hspace*{4pt}$b_\phi=\hat e_\phi-\rho_\phi,\quad \hat y_k=\alpha_k\,b_\phi,\quad \tilde P_k=P^*_k+\hat y_k$\\[4pt]
      \kt{cT!85!black}{Restricciones (todo $t$, $\phi$, $k\in\mathcal N_\phi$)}\\[2pt]
      \textbf{(0)} compra: $\ g\ge\hat P_{L,\mathrm{tot}}-\sum_k\tilde P_k,\quad g\ge0$\\[3pt]
      \textbf{(a)} residual: $\ |\rho_\phi|\le\bar\rho,\quad \sum_\phi\rho_\phi=0$\\[3pt]
      \textbf{(b)} potencia: $\ P^{\min}_k+\alpha_kR^{\varepsilon\downarrow}_\phi\le\tilde P_k\le P^{\max}_k-\alpha_kR^{\varepsilon\uparrow}_\phi$\\[3pt]
      \textbf{(c)} energía: $\ W^{\min}_k(t)+\alpha_kZ^{\downarrow}_\phi(t)\le\sum_{s\le t}\tilde P_k(s)\Delta t$\\[1pt]
      \hspace*{1.55cm}$\sum_{s\le t}\tilde P_k(s)\Delta t\le W^{\max}_k(t)-\alpha_kZ^{\uparrow}_\phi(t)$\\[3pt]
      \textbf{(d)} margen de la fase: $\ \sum_k r_k=R^{\varepsilon}_\phi$ (se cumple solo)\\[3pt]
      \textbf{(e)} final: $\ \sum_{s<T}\tilde P_k(s)\Delta t\le W^{\mathrm{fin}}_k$\\[4pt]
      {\color{grisM}Todo es lineal porque $\alpha_k$, $R^{\varepsilon}$ y $Z$ son datos.}};\end{tikzpicture}
  \end{columns}
\end{frame}

% =====================================================================
%  L7 — LAZO RÁPIDO
% =====================================================================
\begin{frame}{Lazo rápido coherente con el terciario}
  \vspace{-0.30cm}
  \centering
  \resizebox{\linewidth}{!}{%
  \begin{tikzpicture}[x=1cm,y=1cm,font=\small]
    \fill[cL!7,rounded corners=4pt] (2.25,-0.65) rectangle (12.45,3.95);
    \node[font=\small\bfseries,cL!80!black,anchor=south west] at (2.35,-0.6) {Lazo rápido de la fase $\phi$ (en la cabecera, cada $\sim$1 s)};
    \fill[cP!7,rounded corners=4pt] (12.65,-0.65) rectangle (17.45,3.95);
    \node[font=\small\bfseries,cP,anchor=south west] at (12.75,-0.6) {En cada IC de la fase};
    \node[blq=gray!60!black,text width=1.6cm] (M) at (0.95,2.0) {Cabecera\\ mide $P_{h,\phi}$\\ en las 3 fases};
    \node[blq=cL,text width=2.4cm] (EB) at (3.75,2.0) {error de la fase\\ $e_\phi=P_{h,\phi}-\bar P_h$};
    \node[sm=cL] (S) at (5.75,2.0) {$\Sigma$};
    \node[font=\scriptsize] at (5.32,2.38) {$+$};
    \node[font=\scriptsize] at (6.15,1.6) {$-$};
    \node[blq=cL,text width=3.3cm] (PI) at (8.8,2.0) {PI con anti-windup\\ salida limitada a\\ $[\,b_\phi-R^{\varepsilon\downarrow}_\phi,\;b_\phi+R^{\varepsilon\uparrow}_\phi\,]$};
    \node[blq=cL!70!black,text width=2.0cm,font=\scriptsize] (U) at (11.35,0.55) {uso del margen\\ $u_\phi=\dfrac{\Delta P_\phi-b_\phi}{R^{\varepsilon}_\phi}$};
    \draw[ar=gray!60!black] (M) -- (EB);
    \draw[ar=cL] (EB) -- (S);
    \draw[ar=cL] (S) -- node[above,font=\scriptsize]{$e_\phi-\rho_\phi$} (PI);
    \draw[ar=cT,densely dotted] (5.75,0.75) node[below,font=\scriptsize,cT!85!black]{$\rho_\phi$} -- (S.south);
    \draw[ar=cT,densely dotted] (8.8,4.55) node[above,font=\scriptsize,cT!85!black]{$b_\phi,\ R^{\varepsilon}_\phi$ (del terciario)} -- (PI.north);
    \draw[line width=1.3pt,cL] (PI.east) -- node[pos=0.3,above,font=\scriptsize]{$\Delta P_\phi$} (12.95,2.0);
    \draw[ar=cL,densely dotted] (11.35,2.0) -- (U.north);
    \node[jd=cL] at (11.35,2.0) {};
    \draw[line width=1.3pt,cL] (12.95,0.5) -- (12.95,3.4);
    \node[jd=cL] at (12.95,2.0) {};
    \foreach \y/\k in {3.4/1,2.0/2,0.5/N}{
      \node[blq=cP,text width=0.95cm,minimum height=0.65cm] (X\k) at (14.0,\y) {$\times\,\alpha_{\k}$};
      \node[sm=cP] (P\k) at (15.35,\y) {$+$};
      \node[blq=cP,text width=0.8cm,minimum height=0.65cm] (I\k) at (16.75,\y) {IC$_{\k}$};
      \draw[ar=cL] (12.95,\y) -- (X\k.west);
      \draw[ar=cP] (X\k) -- (P\k);
      \draw[ar=cP] (P\k) -- (I\k);
    }
    \draw[ar=cT] (15.35,4.3) node[above,font=\scriptsize,cT!85!black]{$P^*_1$} -- (P1.north);
    \draw[ar=cT] (15.95,2.75) node[right,font=\scriptsize,cT!85!black]{$P^*_2$} -- (P2.north east);
    \draw[ar=cT] (15.95,1.25) node[right,font=\scriptsize,cT!85!black]{$P^*_N$} -- (PN.north east);
    \node[font=\small,cP] at (14.0,1.3) {$\vdots$};
  \end{tikzpicture}}

  \vspace{4pt}
  \begin{columns}[T,onlytextwidth]
    \column{.325\textwidth}
    \begin{tikzpicture}\node[kc=cL,text width=3.8cm,minimum height=1.6cm]{%
      \kt{verdeL}{1. Banda $b_\phi\pm R^{\varepsilon}_\phi$}\\[2pt]
      Ej.: $b=1{,}5$, $R=1{,}37$ kW $\Rightarrow$ $\Delta P\in[0{,}13;\,2{,}87]$ kW. Con $\pm1{,}37$ el PI quedaría saturado toda la tarde solo por lo previsible.};\end{tikzpicture}
    \column{.335\textwidth}
    \begin{tikzpicture}\node[kc=cT,text width=3.9cm,minimum height=1.6cm]{%
      \kt{cT!85!black}{2. La compra no cambia}\\[2pt]
      $\sum_\phi e_\phi=0$ y $\sum_\phi\rho_\phi=0$ $\Rightarrow$ $\sum_\phi\Delta P_\phi=0$ si ninguna fase se satura. El lazo no gasta margen en el residual aceptado.};\end{tikzpicture}
    \column{.325\textwidth}
    \begin{tikzpicture}\node[kc=cP,text width=3.8cm,minimum height=1.6cm]{%
      \kt{cP}{3. Lo mismo, visto en el IC}\\[2pt]
      $P_{\mathrm{ref},k}=P^*_k+\alpha_k\Delta P_\phi$\\
      $\phantom{P_{\mathrm{ref},k}}=P^*_k+\hat y_k+u_\phi\,r_k$\\[2pt]
      $\hat y_k$ puede ir como prealimentación.};\end{tikzpicture}
  \end{columns}
\end{frame}

% =====================================================================
%  L8 — DECISIONES A VALIDAR
% =====================================================================
\begin{frame}{Decisiones que necesito validar}
  \vspace{-0.30cm}
  \begin{tikzpicture}\node[kc=cS,text width=12.25cm,font=\fontsize{7.2}{9}\selectfont]{%
    \renewcommand{\arraystretch}{1.25}%
    \begin{tabular}{@{}>{\bfseries}p{2.9cm}p{4.8cm}p{3.9cm}@{}}
      \textcolor{usachTd}{Decisión} & \textcolor{usachTd}{\textbf{Propuesta}} & \textcolor{usachTd}{\textbf{Alternativa}}\\\hline
      1. Reparto en la fase & $\alpha_k=E_k/\sum_jE_j$ (mismo $\Delta$SoC), fijo en cada resolución & optimizar $\alpha_k$ (LP secuencial)\\
      2. Riesgo y margen & $\varepsilon=1\,\%$, cuantil empírico del error $+\,m$ & gaussiano $z\sigma_\phi+m$\\
      3. Margen de energía & $Z_\phi(t)$ creciente; comparar con solo el primer tramo & solo el primer tramo\\
      4. Lo previsible en el IC & $\hat y_k$ como prealimentación; el PI corrige solo el error & que lo alcance el integrador\\
      5. Residual $\bar\rho$ & límite de desbalance aceptado (norma o diseño) & $\bar\rho=0$ con holgura penalizada\\
      6. Fin del horizonte & SoC objetivo al final de las 12 h ($W^{\mathrm{fin}}_k$) & valorar la energía final en el objetivo\\
      7. Pronóstico & naive estacional (mismo tramo, $n$ días) & corregido con la última medición\\
      8. Lado AC del IC & el IC no regula $V_{dc}$ (lo hacen BESS y secundario) & que el IC regule $V_{dc}$\\
    \end{tabular}};\end{tikzpicture}
\end{frame}

\end{document}
